\documentclass[10pt,a4paper]{amsart}
\usepackage{amsmath,amssymb,amsthm,graphicx}
\usepackage{verbatim}
\newenvironment{reference}{\comment}{\endcomment}
%\newenvironment{reference}{}{}
\newenvironment{slogan}{\comment}{\endcomment}
\newenvironment{history}{\comment}{\endcomment}
% For commutative diagrams we use Xy-pic
\usepackage[all]{xy}

% We use 2cell for 2-commutative diagrams.
\xyoption{2cell}
% Theorem environments.
%
\theoremstyle{plain}
\newtheorem{theorem}[subsection]{Theorem}
\newtheorem{proposition}[subsection]{Proposition}
\newtheorem{lemma}[subsection]{Lemma}

\theoremstyle{definition}
\newtheorem{definition}[subsection]{Definition}
\newtheorem{example}[subsection]{Example}
\newtheorem{exercise}[subsection]{Exercise}
\newtheorem{situation}[subsection]{Situation}

\theoremstyle{remark}
\newtheorem{remark}[subsection]{Remark}
\newtheorem{remarks}[subsection]{Remarks}

\numberwithin{equation}{subsection}

% Macros
%
\def\lim{\mathop{\mathrm{lim}}\nolimits}
\def\colim{\mathop{\mathrm{colim}}\nolimits}
\def\Spec{\mathop{\mathrm{Spec}}}
\def\Hom{\mathop{\mathrm{Hom}}\nolimits}
\def\Ext{\mathop{\mathrm{Ext}}\nolimits}
\def\SheafHom{\mathop{\mathcal{H}\!\mathit{om}}\nolimits}
\def\SheafExt{\mathop{\mathcal{E}\!\mathit{xt}}\nolimits}
\def\Sch{\mathit{Sch}}
\def\Mor{\mathop{\mathrm{Mor}}\nolimits}
\def\Ob{\mathop{\mathrm{Ob}}\nolimits}
\def\Sh{\mathop{\mathit{Sh}}\nolimits}
\def\NL{\mathop{N\!L}\nolimits}
\def\CH{\mathop{\mathrm{CH}}\nolimits}
\def\proetale{{pro\text{-}\acute{e}tale}}
\def\etale{{\acute{e}tale}}
\def\QCoh{\mathit{QCoh}}
\def\Ker{\mathop{\mathrm{Ker}}}
\def\Im{\mathop{\mathrm{Im}}}
\def\Coker{\mathop{\mathrm{Coker}}}
\def\Coim{\mathop{\mathrm{Coim}}}

% Boxtimes
%
\DeclareMathSymbol{\boxtimes}{\mathbin}{AMSa}{"02}

%
% Macros for moduli stacks/spaces
%
\def\QCohstack{\mathcal{QC}\!\mathit{oh}}
\def\Cohstack{\mathcal{C}\!\mathit{oh}}
\def\Spacesstack{\mathcal{S}\!\mathit{paces}}
\def\Quotfunctor{\mathrm{Quot}}
\def\Hilbfunctor{\mathrm{Hilb}}
\def\Curvesstack{\mathcal{C}\!\mathit{urves}}
\def\Polarizedstack{\mathcal{P}\!\mathit{olarized}}
\def\Complexesstack{\mathcal{C}\!\mathit{omplexes}}
% \Pic is the operator that assigns to X its picard group, usage \Pic(X)
% \Picardstack_{X/B} denotes the Picard stack of X over B
% \Picardfunctor_{X/B} denotes the Picard functor of X over B
\def\Pic{\mathop{\mathrm{Pic}}\nolimits}
\def\Picardstack{\mathcal{P}\!\mathit{ic}}
\def\Picardfunctor{\mathrm{Pic}}
\def\Deformationcategory{\mathcal{D}\!\mathit{ef}}
\usepackage{cleveref}
\crefname{originallemma}{original lemma}{original lemmas}
\Crefname{originallemma}{Original lemma}{Original lemmas}
\makeatletter
\newlabel{lemma-check-universally-injective-covering}{{lemma-check-universally-injective-covering}{source}{}{sourcepointer-1}{}}
\newlabel{lemma-check-universally-injective-covering@cref}{{[originallemma][0][]lemma-check-universally-injective-covering}{[1][0][]source}{}{}{}}
\makeatother
\title{Being a universal homeomorphism can be checked after a flat cover of the target}
\author{The Stacks Project Authors}
\date{stacks-project:stacks-morphisms.tex:local-snapshot-20261012; GNU FDL 1.2 or later}
\begin{document}
\maketitle
\noindent\textit{Editorial note.} All delivered mathematical text is the author's own, copied from the pinned Stacks Project chapter source sources/stacks-algebraic-stacks-morphisms/pinned-original-source.tex; no statement or proof step was rewritten, reordered or removed.\par\medskip
\noindent\textit{Editorial note.} Preamble provenance: the Stacks Project keeps one single shared preamble for every chapter; this card takes it directly from the staged sources/stacks-algebraic-stacks-morphisms/preamble.tex as three verbatim spans (lines 8-12, 14-18 and 154-217), so no derived file was written and nothing was added or redefined. The wrapper class and the base packages are supplied by the assembler.\par\medskip
\noindent\textit{Editorial note.} Bibliography provenance: neither the counted statement nor its proof contains a citation (this lane rescans the whole span before writing the card), so the card carries no bibliography and no entry was reconstructed.\par\medskip
\noindent\textit{Editorial note.} Selection provenance: this card was opened under the file-based control contract (CONTROL.md), with the mechanical pre-card gate run on its statement and proof interval before the card was written.\par\medskip
\noindent\textit{Allowed standard prerequisites (editorial).}\begin{itemize}
\item The reduction of universal injectivity to a cover of the target by an algebraic space. Invoked at the start of the counted proof (source lines 3185-3187) as a pointer to a lemma of this chapter. Transfers universal injectivity from the base change to the morphism itself.
\item The openness on topological spaces of surjective flat morphisms of finite presentation, and the formation of the topological space of an algebraic stack as a quotient of the space of a smooth cover. Used inside the counted proof (source lines 3190-3211). Turns the homeomorphism on the cover into openness of the induced map on the topological spaces of the test stacks.
\end{itemize}\medskip
\ifdefined\hypertarget\hypertarget{sourcepointer-1}{}\fi
\noindent\textit{Source reference.} Original lemma lemma-check-universally-injective-covering: sources/stacks-algebraic-stacks-morphisms/pinned-original-source.tex, source lines 3187--3187. Navigation pointer only: the counted text refers to the Stacks item lemma-check-universally-injective-covering, which belongs to another chapter of the project and is not part of this card. The pointer carries the target's own label as its number because this chapter's staged snapshot does not contain the other chapter's tag table; its content is not reproduced here.\par\medskip
\begin{lemma}
\label{lemma-check-universal-homeomorphism-covering}
Let $f : \mathcal{X} \to \mathcal{Y}$ be a morphism of algebraic stacks.
Let $W \to \mathcal{Y}$ be surjective, flat, and locally of finite
presentation where $W$ is an algebraic space. If the base change
$W \times_\mathcal{Y} \mathcal{X} \to W$ is a universal homeomorphism,
then $f$ is a universal homeomorphism.
\end{lemma}

\begin{proof}
Assume $g : W \times_\mathcal{Y} \mathcal{X} \to W$ is a universal
homeomorphism. Then $g$ is universally injective, hence $f$ is
universally injective by
Lemma \ref{lemma-check-universally-injective-covering}.
On the other hand, let $\mathcal{Z} \to \mathcal{Y}$
be a morphism with $\mathcal{Z}$ an algebraic stack.
Choose a scheme $U$ and a surjective
smooth morphism $U \to W \times_\mathcal{Y} \mathcal{Z}$.
Consider the diagram
$$
\xymatrix{
W \times_\mathcal{Y} \mathcal{X} \ar[d]^g &
U \times_\mathcal{Y} \mathcal{X} \ar[d] \ar[l] \ar[r] &
\mathcal{Z} \times_\mathcal{Y} \mathcal{X} \ar[d] \\
W &
U \ar[l] \ar[r] &
\mathcal{Z}
}
$$
The middle vertical arrow induces a homeomorphism
on topological space by assumption on $g$.
The morphism $U \to \mathcal{Z}$ and
$U \times_\mathcal{Y} \mathcal{X} \to
\mathcal{Z} \times_\mathcal{Y} \mathcal{X}$
are surjective, flat, and locally of finite presentation
hence induce open maps on topological spaces.
We conclude that
$|\mathcal{Z} \times_\mathcal{Y} \mathcal{X}| \to |\mathcal{Z}|$
is open. Surjectivity is easy to prove; we omit the proof.
\end{proof}

\end{document}
